\section{Related Work}

\subsection{Alignment Evaluation Benchmarks}

Current alignment evaluation infrastructure focuses predominantly on AI-to-human alignment. RewardBench~\citep{Lambert2024} evaluates reward models on preference prediction across chat, safety, and reasoning domains. AlignBench provides Chinese-language alignment evaluation with fine-grained capability assessment. PERSONA~\citep{PERSONA2024} introduces pluralistic alignment by testing whether models can serve diverse user profiles. These benchmarks share a common assumption: alignment success means the AI accurately predicts and satisfies human preferences.

What these benchmarks do not measure is whether AI responses support or undermine human agency. A response that perfectly matches stated preferences may still lead to cognitive offloading, reduced critical evaluation, or inappropriate reliance.

\subsection{Human Agency in AI Systems}

The concept of agency-preserving AI has emerged from both AI safety and human-computer interaction research. \citet{Mitelut2023} provide a formal framework for understanding how intent-aligned AI systems can deplete human agency by removing decision-making friction that serves epistemic purposes.

HumanAgencyBench~\citep{HumanAgencyBench2024} operationalizes this concern with six measurable dimensions: clarifying ambiguity, presenting options, epistemic hedging, explicit deferral, supporting user reasoning, and avoiding cognitive shortcuts. We adapt four of these dimensions as computable proxies for our BAI computation.

\subsection{Bidirectional Alignment Framework}

\citet{Shen2024} synthesize concerns about unidirectional alignment into a Bidirectional Alignment Framework. Their systematic review of 400+ papers demonstrates that the field overwhelmingly focuses on AI$\rightarrow$Human alignment while neglecting Human$\rightarrow$AI alignment dimensions. Our work directly addresses this operationalization gap.

\subsection{Adversarial Probing}

Gradient reversal training, introduced by \citet{Ganin2015} for domain adaptation, provides a method for isolating orthogonal representational components. We apply gradient reversal to investigate whether BAI occupies a representationally independent subspace from reward prediction.
